\pdfoutput=1
\pdfinfoomitdate=1
\pdftrailerid{}
\pdfsuppressptexinfo=15
\documentclass[indonesian]{shinybook}
\usepackage{tikz}
\usepackage{pgfplots}
\usepgfplotslibrary{groupplots}
\pgfplotsset{compat=1.18}
\usetikzlibrary{intersections,calc}
\usetikzlibrary{arrows.meta, decorations.markings}
\usepackage{glossaries-extra}
\usepackage{needspace}

\newcommand{\half}{\frac{1}{2}}
\newcommand{\N}{\mathbb{N}}
\newcommand{\Z}{\mathbb{Z}}
\newcommand{\Q}{\mathbb{Q}}
\newcommand{\R}{\mathbb{R}}
\newcommand{\Rb}{\overline{\mathbb{R}}}
\newcommand{\bR}{\overline{\mathbb{R}}}
\newcommand{\F}{\mathbb{F}}
\newcommand{\1}{\mathbb{1}}
\newcommand{\linop}{\mathcal{L}}
\newcommand{\Rnn}{\mathbb{R}^{n\times n}}
\newcommand{\Cnn}{\mathbb{C}^{n\times n}}
\newcommand{\E}{\mathbb{E}}
\usepackage{pgffor}
\foreach \x in {A,...,Z}{%
    \expandafter\xdef\csname cal\x\endcsname{\noexpand\ensuremath{\noexpand\mathcal{\x}}}
}

\renewcommand{\implies}{\Rightarrow}
\newcommand{\setto}{\rightrightarrows}
\newcommand{\equivalent}{\Leftrightarrow}
\newcommand{\eps}{\varepsilon}
\renewcommand{\phi}{\varphi}
\newcommand{\supp}{\operatorname{\mathrm{supp}}}
\renewcommand{\div}{\operatorname{\mathrm{div}}}
\newcommand{\dom}{\operatorname{\mathrm{dom}}}
\renewcommand{\ker}{\operatorname{\mathrm{ker}}}
\newcommand{\conv}{\operatorname{\mathrm{conv}}}
\newcommand{\rg}{\operatorname{\mathrm{rg}}}
\newcommand{\graph}{\operatorname{\mathrm{graph}}}
\newcommand{\tr}{\operatorname{\mathrm{tr}}}
\newcommand{\epi}{\operatorname{\mathrm{epi}}}
% Never redefine TeX's primitive \span: amsmath alignment environments use it.
% Use \spann for the linear-span operator in reader prose and formulae.
\newcommand{\sign}{\operatorname{\mathrm{sign}}}
\DeclareMathOperator{\Id}{\mathrm{Id}}
\newcommand{\prox}{\mathrm{prox}}
\newcommand{\proj}{\mathrm{proj}}
\newcommand{\spann}{\operatorname{\mathrm{span}}}
\newcommand{\Lloc}{L^1_{\mathrm{loc}}}
\newcommand{\Cinf}{C^\infty_0}

% Nama lingkungan dilokalkan; penomoran dan topologi sumber dipertahankan.
\newtheorem{theorem}{Teorema}[chapter]
\newtheorem{lemma}[theorem]{Lema}
\newtheorem{cor}[theorem]{Korolari}
\theoremstyle{definition}
\newtheorem{defn}[theorem]{Definisi}
\newtheorem{prop}[theorem]{Proposisi}
\newtheorem{example}[theorem]{Contoh}
\newtheorem{assumption}[theorem]{Asumsi}
\newtheorem{exercise}[theorem]{Latihan}
\newtheorem*{rem}{Catatan}

\crefname{theorem}{teorema}{teorema}
\Crefname{theorem}{Teorema}{Teorema}
\crefname{lemma}{lema}{lema}
\Crefname{lemma}{Lema}{Lema}
\crefname{cor}{korolari}{korolari}
\Crefname{cor}{Korolari}{Korolari}
\crefname{prop}{proposisi}{proposisi}
\Crefname{prop}{Proposisi}{Proposisi}
\crefname{defn}{definisi}{definisi}
\Crefname{defn}{Definisi}{Definisi}
\crefname{example}{contoh}{contoh}
\Crefname{example}{Contoh}{Contoh}
\crefname{assumption}{asumsi}{asumsi}
\Crefname{assumption}{Asumsi}{Asumsi}
\crefname{rem}{catatan}{catatan}
\Crefname{rem}{Catatan}{Catatan}
\crefname{exercise}{latihan}{latihan}
\Crefname{exercise}{Latihan}{Latihan}

\renewcommand{\proofname}{Bukti}
\renewcommand{\contentsname}{Daftar Isi}
\renewcommand{\figurename}{Gambar}

\newcommand{\norm}[1]{\|#1\|}
\newcommand{\abs}[1]{\left|#1\right|}
\newcommand{\inner}[2]{\left\langle#1 , #2\right\rangle}
\newcommand{\dual}[1]{\langle #1 \rangle}
\newcommand{\setof}[2]{\left\{#1 \;\middle|\; #2\right\}}
\newcommand{\wkto}{\rightharpoonup}
\DeclareMathOperator{\Exp}{\mathbb{E}}
\DeclareMathOperator{\Var}{\mathbb{V}}
\DeclareMathOperator{\Cov}{\mathrm{Cov}}
\newcommand{\Prob}{\mathbb{P}}

\usepackage{enumerate}
\newcommand{\weaki}{\protect{\u{\i}}}
\AtBeginEnvironment{remark}{\small}

\ifdefined\OIntegratedReader
  % Keep examples as ordinary theorem blocks in the tagged master. The legacy
  % mdframed/tablefootnote patches interfere with the synchronized latex-lab
  % hooks, while no canonical body calls \tablefootnote.
\else
  \usepackage{mdframed}
  \mdfdefinestyle{examplebg}{
      backgroundcolor=structure!7,%
      hidealllines=true,%
      splittopskip=\topskip,
      frametitleaboveskip=0,
      frametitlebelowskip=0,
      innertopmargin=-.3em,%
      innerbottommargin=.7em,%
      innerleftmargin=.7em,%
      innerrightmargin=.7em,%
  }
  \surroundwithmdframed[style=examplebg]{example}
  \usepackage{tablefootnote}
  \makeatletter
  \AfterEndEnvironment{mdframed}{%
   \tfn@tablefootnoteprintout%
   \gdef\tfn@fnt{0}%
  }
  \makeatother
\fi

\usepackage{enumitem}
\setlist[enumerate]{label=(\roman*)}
\usepackage{scrhack}

\newcommand{\ie}{\textit{yaitu}}
\newcommand{\eg}{\textit{misalnya}}
\newcommand{\cf}{\textit{bandingkan}}
\newcommand{\emptyarg}{\,\cdot\,}
\newcommand{\dd}{\mathrm{d}}
\newcommand{\todo}[1]{{\color{red}Perlu dikerjakan: #1}}
\newcommand{\Oc}{\mathcal{O}}
\newcommand{\Lc}{\mathcal{L}}
\newcommand{\Gc}{\mathcal{G}}
\newcommand{\Ic}{\mathcal{I}}
\newcommand{\Pc}{\mathcal{P}}
\newcommand{\Mc}{\mathcal{M}}
\renewcommand{\epsilon}{\varepsilon}
\newcommand{\diag}{\mathrm{diag}}

% Salinan metadata lokal ini memperbaiki nama pengarang tunggal pada entri
% Villani tanpa mengubah berkas otoritas yang dibekukan.
\addbibresource{references-ot-id.bib}
\microtypesetup{expansion=false}
\glsdisablehyper
\setlength{\emergencystretch}{2em}

% Lokalisasi ini sengaja hanya berlaku pada unit mandiri Bab 9.
\crefname{equation}{persamaan}{persamaan}
\Crefname{equation}{Persamaan}{Persamaan}

\title{Optimisasi Konveks}
\subtitle{Unit 9: Transportasi Optimal\\Edisi Bahasa Indonesia}
\author{Andreas Habring}
\date{Sumber edisi v1: 13 Juli 2026}
\publishers{\small
Terjemahan kerja bahasa Indonesia dari \emph{Lecture Notes: Convex Optimization}, arXiv:2607.11664v1.\\
Sumber asli dan terjemahan ini dilisensikan CC BY 4.0. Perubahan matematis dijelaskan dalam catatan koreksi.\\
Terjemahan mandiri ini bukan karya resmi atau dukungan Andreas Habring maupun TU Graz.}
\hypersetup{
    pdftitle={Optimisasi Konveks - Unit 9: Transportasi Optimal - Edisi Bahasa Indonesia},
    pdfauthor={Andreas Habring},
    pdflang={id-ID},
}

\begin{document}
\maketitle

\frontmatter
\chapter*{Tentang unit ini}
Unit ini menerjemahkan secara utuh Bab 9, "Excursion on Optimal Transport", dari Andreas Habring, \emph{Lecture Notes: Convex Optimization}, arXiv:2607.11664v1. Sumber berwenang untuk unit ini adalah berkas \texttt{optimal\_transport.tex} edisi v1 dengan SHA-256 berikut:
\begin{center}
\small\texttt{719df724b368126cc7540dffd461dc33}\\[-0.2em]
\texttt{aba7d5b5b6060132181086dfa17649ba}
\end{center}

Bab-bab prasyarat tidak diulang. Pembaca diharapkan telah mengenal ruang terukur dan ukuran probabilitas, konvergensi lemah, ruang Polish dan keketatan, optimisasi konveks, dualitas, serta aljabar matriks. Seluruh definisi, gambar TikZ, teorema, sketsa bukti, lema, catatan, sitasi, catatan kaki, formula, rujukan silang, dan permukaan glosarium dalam bab sumber dipertahankan dalam urutan sumber. Kekeliruan yang dapat ditentukan secara matematis diperbaiki secara terbuka dan diringkas pada akhir unit.

Uraian sumber menyebut dirinya sebagai selingan yang tidak sepenuhnya rigor. Edisi ini mempertahankan sifat pengantar tersebut, sambil menyatakan hipotesis minimum yang diperlukan agar klaim eksistensi, metrik Wasserstein, dualitas Kantorovich, serta faktorisasi dan iterasi Sinkhorn mempunyai makna yang tepat.

\tableofcontents

% Rujukan mandiri ke Bab 7 dibuat bernomor tanpa memasukkan ulang bab tersebut.
\setcounter{chapter}{6}
\refstepcounter{chapter}\label{chapter:duality}

\mainmatter
\setcounter{chapter}{8}
% H09-S001 | sumber authority/habring/source-v1/optimal_transport.tex baris 1--15
% segment-id: d90.hab.v1.ch09.seg0001
\newabbreviation{ot}{OT}{transportasi optimal}


\chapter{Selingan tentang Transportasi Optimal}
Meskipun tidak sepenuhnya menyatu dengan kuliah ini, kita menutupnya dengan suatu selingan mengenai \gls{ot}, bidang yang sangat penting dalam matematika modern dan pembelajaran mesin. Perlu ditekankan sejak awal bahwa uraian ini hanya pengantar singkat; sejumlah rincian teori ukuran dan topologi yang mendasarinya akan dinyatakan seperlunya.

\paragraph{Notasi}
Dalam bagian berikut, $(X,\Sigma_X)$ dan $(Y,\Sigma_Y)$ adalah ruang terukur\footnote{Ketika topologi dibutuhkan, keduanya akan diasumsikan sebagai ruang Polish dengan aljabar-$\sigma$ Borel; yaitu, ruang topologis separabel yang topologinya dibangkitkan oleh suatu metrik lengkap.}. Tuliskan $\Pc(X)$ dan $\Pc(Y)$ untuk himpunan semua ukuran probabilitas, $\Mc(X)$ dan $\Mc(Y)$ untuk ruang ukuran bertanda hingga, serta $\Mc_+(X)$ dan $\Mc_+(Y)$ untuk kerucut ukuran tak negatif hingga.
Ingat bahwa ukuran bertanda $\mu$ pada $X$ adalah fungsi $\mu:\Sigma_X\rightarrow\R$, sedangkan ukuran tak negatif bernilai dalam $[0,\infty)$, dengan $\Sigma_X\subseteq2^X$ suatu aljabar-$\sigma$\footnote{Secara khusus, $\Sigma_X$ memuat himpunan kosong serta tertutup terhadap komplemen dan gabungan terhitung.}, dan memenuhi
\begin{enumerate}
    \item $\mu(\emptyset)=0$;
    \item untuk setiap barisan $(A_n)_{n\in\N}\subset\Sigma_X$ yang saling lepas, berlaku $\mu(\bigcup_n A_n)=\sum_n\mu(A_n)$.
\end{enumerate}
Ukuran probabilitas adalah ukuran tak negatif dengan $\mu(X)=1$.

% H09-S002 | sumber authority/habring/source-v1/optimal_transport.tex baris 16--75
% segment-id: d90.hab.v1.ch09.seg0002
\section{Transportasi optimal Monge dan Kantorovich}

\begin{figure}
\centering
\resizebox{\textwidth}{!}{%
\begin{tikzpicture}[
  >=Stealth,
  font=\sffamily
]

% ----- garis dasar bersama -----
\def\y{0}
\def\xmin{0}
\def\xmax{16}
\def\amean{3}
\def\astd{1}
\def\bmean{12}
\def\bstd{1.4}

% ----- kurva kepadatan tegak pada garis dasar yang sama -----
% alpha: Gaussian, rerata \amean, simpangan baku \astd
\draw[domain=\xmin:\xmax, smooth, samples=150, variable=\x,
      blue!60!black, thick, fill=blue!15]
  plot ({\x}, {\y + 1.6*exp(-0.5*((\x-\amean)/\astd)^2)})
  -- (\xmax,\y) -- (\xmin,\y) -- cycle;

% beta: Gaussian, rerata \bmean, simpangan baku \bstd
\draw[domain=\xmin:\xmax, smooth, samples=150, variable=\x,
      red!60!black, thick, fill=red!15]
  plot ({\x}, {\y + 1.6*exp(-0.5*((\x-\bmean)/\bstd)^2)})
  -- (\xmax,\y) -- (\xmin,\y) -- cycle;

% ----- sumbu garis dasar bersama -----
\draw[gray] (\xmin,\y) -- (\xmax,\y);

% ----- label distribusi -----
\node[blue!60!black] at (\amean, 2.1) {$\alpha$};
\node[red!60!black]  at (\bmean, 2.1) {$\beta$};

% ----- peta transportasi: busur monoton tak berpotongan di bawah sumbu -----
% titik sampel pada dukungan alpha dipetakan ke dukungan beta oleh
% peta optimal monoton T(x)=\bmean+(\bstd/\astd)(x-\amean)
\foreach \x/\opacity in {1/0.4, 2/0.6, 3/0.9, 4/0.6, 5/0.4} {
  \pgfmathsetmacro{\Tx}{\bmean + (\bstd/\astd)*(\x-\amean)}
  \draw[-{Stealth[length=2.2mm]}, line width=1pt,
        violet!70!black, opacity=\opacity]
    (\x,\y) to[out=-60, in=-120] (\Tx,\y);
  \filldraw[blue!60!black] (\x,\y) circle (1.3pt);
  \filldraw[red!60!black] (\Tx,\y) circle (1.3pt);
}

% ----- label peta -----
\node[violet!70!black, fill=white, fill opacity=0.85, text opacity=1,
      inner sep=2pt] at (8, -3.7) {$T_{\#}\alpha=\beta$};
\node[violet!70!black, font=\sffamily\small, fill=white, fill opacity=0.85,
      text opacity=1, inner sep=2pt] at (8, -4.2) {(peta transportasi optimal)};
% ----- judul opsional -----
% \node[font=\sffamily\bfseries] at (8, 3.2) {Transportasi Optimal dari $\alpha$ ke $\beta$};
\end{tikzpicture}%
}
\caption{Ilustrasi~\gls{ot}. Peta $T$ mengangkut \emph{tumpukan} dari $\alpha$ ke $\beta$.}
\label{ot:fig:ot}
\end{figure}

% H09-S003 | sumber authority/habring/source-v1/optimal_transport.tex baris 77--110
% segment-id: d90.hab.v1.ch09.seg0003
Motivasi khas untuk \gls{ot} adalah sebagai berikut. Bayangkan kita ingin memindahkan setumpuk pasir dari satu tempat ke tempat lain. Berdasarkan kekekalan massa, kedua tumpukan mempunyai massa total yang sama; tanpa mengurangi keumuman, kita andaikan massanya satu. Karena itu, tumpukan awal dapat digambarkan oleh distribusi probabilitas $\alpha\in\Pc(X)$ dan tumpukan akhir oleh distribusi probabilitas $\beta\in\Pc(Y)$. Pertanyaan yang hendak dijawab dalam \gls{ot} adalah:
\begin{quote}
    \centering
    \emph{Bagaimana kita dapat mengangkut tumpukan $\alpha$ ke $\beta$ dengan biaya minimum?}
\end{quote}
Pembaca yang teliti mungkin menyadari bahwa \emph{biaya} belum ditentukan. Andaikan tersedia fungsi biaya terukur yang terbatas dari bawah, $c:X\times Y\rightarrow(-\infty,+\infty]$, dengan $c(x,y)$ menyatakan biaya untuk mengangkut satu satuan massa dari $x\in X$ ke $y\in Y$. Setiap transportasi deterministik dari $\alpha$ ke $\beta$ dimodelkan oleh pemetaan terukur $T:X\rightarrow Y$. Syarat pengangkutan dapat dirumuskan dengan meminta $\alpha(T^{-1}(A))=\beta(A)$ untuk setiap $A\in\Sigma_Y$. Ini mengantar kita pada definisi berikut.
\begin{defn}[Ukuran hasil dorong (\emph{push-forward})]
    Misalkan $\alpha\in\Pc(X)$ dan $T:X\rightarrow Y$ terukur. Ukuran hasil dorong $T_\sharp\alpha\in\Pc(Y)$ didefinisikan oleh $T_\sharp\alpha(A)=\alpha(T^{-1}(A))$ untuk setiap $A\in\Sigma_Y$.
\end{defn}

Sekarang kita dapat memperkenalkan secara formal masalah Monge untuk~\gls{ot}:
\begin{defn}[Monge \gls{ot}]
    Misalkan $X$ dan $Y$ ruang terukur, $\alpha\in\Pc(X)$, $\beta\in\Pc(Y)$, dan $c$ fungsi biaya terukur yang terbatas dari bawah. Masalah Monge untuk \gls{ot} adalah
    \begin{equation}\label{ot:eq:monge}
        \inf_{\substack{T:X\rightarrow Y\ \mathrm{terukur}\\T_\sharp\alpha=\beta}}
        \int c(x,T(x))\dd\alpha(x),
        \qquad \inf\emptyset\coloneqq+\infty.
    \end{equation}
\end{defn}
Masalah Monge untuk \gls{ot} mempunyai beberapa kekurangan. Sebagai contoh, setiap $x$ hanya dapat dipetakan oleh $T$ ke \emph{satu} titik $y$, sehingga massa pada satu titik $x$ tidak dapat dibagi ke beberapa titik $y$. Karena itu, apabila $\alpha$ memberikan massa tak nol pada suatu titik, $Y\subseteq\R^q$, dan $\beta$ kontinu mutlak terhadap ukuran Lebesgue\footnote{Hal ini khususnya menyiratkan bahwa $\beta$ tidak mempunyai massa titik.}, tidak ada peta $T$ yang memenuhi syarat. Selain itu, bergantung pada pilihan $c$, masalah ini umumnya sulit dianalisis karena peubah tak diketahui $T$ muncul di dalam biaya $c$.

Perumusan yang lebih umum dan lebih mudah dianalisis adalah \gls{ot} Kantorovich. Alih-alih memakai peta $T:X\rightarrow Y$, transportasi juga dimodelkan oleh suatu distribusi. Suatu rencana transportasi adalah distribusi probabilitas $\gamma$ pada ruang hasil kali terukur $X\times Y$ yang mempunyai $\alpha$ dan $\beta$ sebagai marginalnya: $\gamma(A\times Y)=\alpha(A)$ dan $\gamma(X\times B)=\beta(B)$ untuk semua $A\in\Sigma_X$ dan $B\in\Sigma_Y$. Secara intuitif, $\gamma(A\times B)$ adalah massa yang diangkut dari $A$ ke $B$. Himpunan semua distribusi dengan marginal $\alpha,\beta$ dilambangkan dengan $\Pi(\alpha,\beta)$ dan anggotanya disebut kopling (\emph{coupling}). Sekarang kita dapat mendefinisikan masalah Kantorovich untuk~\gls{ot}.
\begin{defn}[Kantorovich \gls{ot}]
    Misalkan $X$ dan $Y$ ruang terukur, $\alpha\in\Pc(X)$, $\beta\in\Pc(Y)$, dan $c$ fungsi biaya terukur. Masalah Kantorovich untuk \gls{ot} adalah
    \begin{equation}\label{ot:eq:K_ot}
        \inf_{\gamma\in\Pi(\alpha,\beta)}
        \int c(x,y)\dd\gamma(x,y).
    \end{equation}
\end{defn}
Perhatikan bahwa~\eqref{ot:eq:K_ot} jauh lebih menguntungkan daripada~\eqref{ot:eq:monge}. Secara khusus, objektifnya kini linear.

\begin{rem}
    Kita juga dapat merumuskan~\eqref{ot:eq:K_ot} dengan peubah acak. Setiap kopling $\gamma$ dapat \emph{direalisasikan} sebagai vektor acak $(X,Y)\sim\gamma$ dengan $X\sim\alpha$ dan $Y\sim\beta$. Biayanya kemudian memenuhi
    \begin{equation}
        \int c(x,y)\dd\gamma(x,y)=\E[c(X,Y)].
    \end{equation}
\end{rem}

% H09-S004 | sumber authority/habring/source-v1/optimal_transport.tex baris 112--130
% segment-id: d90.hab.v1.ch09.seg0004
Kita nyatakan hasil berikut.
\begin{theorem}[Eksistensi solusi]
    Misalkan $X$ dan $Y$ ruang Polish yang dilengkapi aljabar-$\sigma$ Borelnya, $\alpha\in\Pc(X)$, dan $\beta\in\Pc(Y)$. Jika $c:X\times Y\rightarrow(-\infty,+\infty]$ terbatas dari bawah dan semikontinu bawah, maka terdapat rencana transportasi optimal $\hat\gamma\in\Pi(\alpha,\beta)$, yaitu solusi untuk~\eqref{ot:eq:K_ot}.
\end{theorem}
\begin{proof}[Sketsa bukti]
    Untuk bukti terperinci, lihat misalnya~\cite[Theorem 4.1]{villani2009optimal}. Argumennya adalah metode langsung. Ambil barisan $(\gamma_n)_n\subset\Pi(\alpha,\beta)$ yang nilai biayanya menuju infimum. Karena $\alpha$ dan $\beta$ ketat pada ruang Polish, keluarga kopling dengan kedua marginal tetap juga ketat. Himpunan $\Pi(\alpha,\beta)$ tertutup terhadap konvergensi lemah; teorema Prokhorov karena itu menyatakan bahwa himpunan tersebut kompak secara lemah. Setelah mengambil suatu subbarisan, diperoleh $\gamma_n\rightharpoonup\hat\gamma\in\Pi(\alpha,\beta)$. Karena $c$ semikontinu bawah dan terbatas dari bawah, pemetaan $\gamma\mapsto\int c(x,y)\dd\gamma(x,y)$ semikontinu bawah secara lemah. Maka $\hat\gamma$ meminimumkan biaya dan bukti selesai.
\end{proof}

Satu kasus khusus penting dari transportasi optimal perlu disoroti.

\begin{defn}[Jarak Wasserstein]
Misalkan $(X,d)$ ruang metrik Polish dan $p\in[1,\infty)$. Tetapkan $x_0\in X$ dan tuliskan
$\Pc_p(X)\coloneqq\{\mu\in\Pc(X):\int d(x,x_0)^p\dd\mu(x)<\infty\}$; definisi ini tidak bergantung pada pilihan $x_0$. Untuk $\alpha,\beta\in\Pc_p(X)$ dan biaya $c(x,y)=d(x,y)^p$, besaran
\begin{equation}
    W_p(\alpha,\beta)\coloneqq
    \bigg(\inf_{\gamma\in\Pi(\alpha,\beta)}
    \int d(x,y)^p\dd\gamma(x,y)\bigg)^{1/p}
\end{equation}
disebut jarak Wasserstein-$p$.
\end{defn}
Secara khusus, $W_p$ adalah metrik pada ruang ukuran probabilitas bermomen ke-$p$ hingga, yaitu $\Pc_p(X)$.

% H09-S005 | sumber authority/habring/source-v1/optimal_transport.tex baris 131--141
% segment-id: d90.hab.v1.ch09.seg0005
\section{Dualitas Kantorovich}
Dari~\cref{chapter:duality} kita telah mengenal konsep dualitas. Dualitas memainkan peran penting dalam~\gls{ot} dan menghasilkan teorema terkenal berikut.
\begin{theorem}[Dualitas Kantorovich]
    Misalkan $X$ dan $Y$ ruang Polish dengan aljabar-$\sigma$ Borelnya, $\alpha\in\Pc(X)$, $\beta\in\Pc(Y)$, dan $c:X\times Y\rightarrow[0,\infty]$ semikontinu bawah. Tuliskan $C_b(X)$ dan $C_b(Y)$ untuk fungsi kontinu terbatas bernilai riil. Maka
    \begin{equation}
        \begin{aligned}
            \inf_{\gamma\in\Pi(\alpha,\beta)}\int c(x,y)\dd\gamma(x,y)
            = \sup_{\substack{\phi\in C_b(X),\ \psi\in C_b(Y)\\
            \phi(x)+\psi(y)\leq c(x,y)\ \forall(x,y)}}
            \int\phi(x)\dd\alpha(x)+\int\psi(y)\dd\beta(y).
        \end{aligned}
    \end{equation}
\end{theorem}

% H09-S006 | sumber authority/habring/source-v1/optimal_transport.tex baris 142--190
% segment-id: d90.hab.v1.ch09.seg0006
\begin{proof}[Sketsa bukti]
    Kita mulai dengan ketaksamaan yang mudah. Ambil $(\phi,\psi)\in C_b(X)\times C_b(Y)$ sedemikian sehingga $\phi(x)+\psi(y)\leq c(x,y)$ untuk semua $(x,y)$, dan ambil sembarang $\gamma\in\Pi(\alpha,\beta)$. Maka
    \begin{equation}
        \begin{aligned}
            \int\phi(x)\dd\alpha(x)+\int\psi(y)\dd\beta(y)
            &= \int\phi(x)\dd\gamma(x,y)+\int\psi(y)\dd\gamma(x,y)\\
            &= \int\big(\phi(x)+\psi(y)\big)\dd\gamma(x,y)\\
            &\leq \int c(x,y)\dd\gamma(x,y).
        \end{aligned}
    \end{equation}
    Karena $\gamma$ dan $(\phi,\psi)$ sembarang, kita dapat mengambil infimum di ruas kanan dan supremum di ruas kiri untuk memperoleh
    \begin{equation}
        \begin{aligned}
            \sup_{\substack{\phi\in C_b(X),\ \psi\in C_b(Y)\\
            \phi(x)+\psi(y)\leq c(x,y)\ \forall(x,y)}}
            \int\phi(x)\dd\alpha(x)+\int\psi(y)\dd\beta(y)
            \leq \inf_{\gamma\in\Pi(\alpha,\beta)}
            \int c(x,y)\dd\gamma(x,y).
        \end{aligned}
    \end{equation}
    Ketaksamaan sebaliknya lebih rumit.
    Pertama, kita dapat menulis ulang~\eqref{ot:eq:K_ot} sebagai
    \begin{equation}
        \inf_{\gamma\in\Mc_+(X\times Y)}
        \left\{\int c(x,y)\dd\gamma+\delta_{\Pi(\alpha,\beta)}(\gamma)\right\},
    \end{equation}
    dengan $\delta_{\Pi(\alpha,\beta)}$ fungsi indikator seperti biasa. Untuk $\gamma\in\Mc_+(X\times Y)$, berlaku
    \begin{equation}\label{ot:eq:duality}
        \begin{aligned}
            \delta_{\Pi(\alpha,\beta)}(\gamma)
            &= \sup_{\phi\in C_b(X),\ \psi\in C_b(Y)}
            \bigg\{\int\phi(x)\dd\alpha(x)+\int\psi(y)\dd\beta(y)\\
            &\qquad-\int\big(\phi(x)+\psi(y)\big)\dd\gamma(x,y)\bigg\}.
        \end{aligned}
    \end{equation}
    Jika $\gamma\in\Pi(\alpha,\beta)$, ekspresi di dalam supremum pada ruas kanan bernilai nol untuk setiap $(\phi,\psi)$, sehingga supremumnya nol. Sebaliknya, jika $\gamma\notin\Pi(\alpha,\beta)$, sedikitnya satu marginal berbeda. Ukuran Borel hingga yang berbeda pada ruang Polish dapat dipisahkan oleh suatu fungsi kontinu terbatas; tanpa mengurangi keumuman, terdapat $\phi\in C_b(X)$ sedemikian sehingga
    \begin{equation}
        \int\phi\dd(\alpha-\gamma_x)\neq0,
    \end{equation}
    dengan $\gamma_x$ marginal-$x$ dari $\gamma$. Mengganti $\phi$ dengan $t\phi$ dan memilih tanda $t$ ketika $|t|\rightarrow\infty$ menunjukkan bahwa supremumnya adalah $+\infty$.

    Langkah taktrivial untuk ketaksamaan sebaliknya adalah pertukaran infimum dan supremum. Pertukaran itu tidak mengikuti perhitungan formal semata; di bawah hipotesis ruang Polish dan biaya semikontinu bawah tak negatif di atas, pertukaran tersebut dijamin oleh teorema dualitas Kantorovich standar, lihat~\cite[Theorem 5.10]{villani2009optimal}. Dengan justifikasi dualitas kuat ini, diperoleh
    \begin{equation}
        \begin{aligned}
            \inf_{\gamma\in\Pi(\alpha,\beta)}
            \int c(x,y)\dd\gamma(x,y)
            &= \inf_{\gamma\in\Mc_+(X\times Y)}
            \sup_{\phi\in C_b(X),\ \psi\in C_b(Y)}
            \\
            &\quad\bigg\{\int c(x,y)\dd\gamma(x,y)\\
            &\quad+\int\phi(x)\dd\alpha(x)+\int\psi(y)\dd\beta(y)
            \\
            &\quad-\int\big(\phi(x)+\psi(y)\big)\dd\gamma(x,y)\bigg\}\\
            &= \sup_{\phi\in C_b(X),\ \psi\in C_b(Y)}
            \bigg\{\int\phi(x)\dd\alpha(x)+\int\psi(y)\dd\beta(y)\\
            &\quad+\inf_{\gamma\in\Mc_+(X\times Y)}
            \int\big(c(x,y)-\phi(x)-\psi(y)\big)\dd\gamma(x,y)\bigg\}\\
            &= \sup_{\substack{\phi\in C_b(X),\ \psi\in C_b(Y)\\
            \phi(x)+\psi(y)\leq c(x,y)\ \forall(x,y)}}
            \int\phi(x)\dd\alpha(x)+\int\psi(y)\dd\beta(y).
        \end{aligned}
    \end{equation}
    Pada persamaan terakhir, infimum terhadap ukuran tak negatif bernilai nol jika $c-\phi-\psi\geq0$ di setiap titik (ambil $\gamma=0$), dan bernilai $-\infty$ jika ketaksamaan itu gagal (skalakan massa Dirac pada titik pelanggaran). Ini sama dengan fungsi indikator kendala dual, sebagaimana pada~\eqref{ot:eq:duality}.
\end{proof}

% H09-S007 | sumber authority/habring/source-v1/optimal_transport.tex baris 192--226
% segment-id: d90.hab.v1.ch09.seg0007
\section{Diskretisasi transportasi optimal: Algoritme Sinkhorn--Knopp}
Untuk mendiskretisasi~\eqref{ot:eq:K_ot}, kita modelkan $\alpha$ dan $\beta$ sebagai ukuran diskret. Tuliskan
$\Delta_n\coloneqq\{r\in\R_+^n:\1_n^\top r=1\}$ dan
$\Delta_n^\circ\coloneqq\{r\in\R_{++}^n:\1_n^\top r=1\}$. Untuk suatu vektor $a=(a_1,\dots,a_n)\in\Delta_n$ dan titik $x_i\in X$, tetapkan
\begin{equation}
    \alpha=\sum_{i=1}^n a_i\delta_{x_i},
\end{equation}
dan, serupa dengan itu, untuk $b=(b_1,\dots,b_m)\in\Delta_m$ dan $y_j\in Y$, tetapkan
\begin{equation}
    \beta=\sum_{j=1}^m b_j\delta_{y_j}.
\end{equation}
Selanjutnya, definisikan $C=[C_{i,j}]_{i,j}=[c(x_i,y_j)]_{i,j}$. Rencana transportasi adalah matriks $P\in\R^{n\times m}$ sedemikian sehingga
\begin{equation}
    \begin{cases}
        \sum_{j=1}^mP_{i,j}&=a_i,\\
        \sum_{i=1}^nP_{i,j}&=b_j,\\
        P_{i,j}&\geq0.
    \end{cases}
\end{equation}
Kendala marginal dapat ditulis ringkas sebagai $P\1_m=a$ dan $P^\top\1_n=b$. Masalah \gls{ot} diskret kemudian berbentuk
\begin{equation}
    \begin{aligned}
        \min_{P\in\R^{n\times m}}\inner{C}{P}
        &=\sum_{i,j}C_{i,j}P_{i,j}\\
        \text{dengan syarat }&
        P\1_m=a,\quad P^\top\1_n=b,\quad P_{i,j}\geq0.
    \end{aligned}
\end{equation}
Dalam praktik, masalah ini sering diregularisasi. Untuk pembahasan regularisasi entropik berikut, andaikan $C\in\R^{n\times m}$, sehingga semua biaya diskret berhingga. Untuk $P\in\R_+^{n\times m}$, definisikan regularisasi entropik
\begin{equation}
    E(P)\coloneqq\sum_{i,j}P_{i,j}\big(\log(P_{i,j})-1\big),
\end{equation}
dengan konvensi kontinu $0\log0=0$; di luar ortan tak negatif, tetapkan $E(P)=+\infty$. Untuk $\epsilon>0$, pertimbangkan masalah
\begin{equation}\label{ot:eq:disc_entropic}
    \begin{aligned}
        \min_{P\in\R_+^{n\times m}}\inner{C}{P}+\epsilon E(P)
        \quad\text{dengan syarat }P\1_m=a,\quad P^\top\1_n=b.
    \end{aligned}
\end{equation}


% H09-S008 | sumber authority/habring/source-v1/optimal_transport.tex baris 229--256
% segment-id: d90.hab.v1.ch09.seg0008
\begin{lemma}
    Andaikan $a\in\Delta_n^\circ$, $b\in\Delta_m^\circ$, $C\in\R^{n\times m}$, dan $\epsilon>0$. Solusi masalah~\eqref{ot:eq:disc_entropic} unik dan mempunyai representasi
    \begin{equation}
        P_{i,j}=u_iK_{i,j}v_j,
    \end{equation}
    dengan $u_i,v_j>0$ dan $K_{i,j}=\exp\left(-\frac{C_{i,j}}{\epsilon}\right)>0$. Vektor penskalaan $(u,v)$ unik hingga transformasi $(u,v)\mapsto(tu,t^{-1}v)$ untuk $t>0$.
\end{lemma}
\begin{proof}
    Himpunan layak adalah politop transportasi tak kosong dan kompak: matriks
    $Q=ab^\top$ layak, bahkan $Q_{i,j}>0$. Karena fungsi objektif kontinu pada
    politop ini, suatu peminimum ada. Fungsi
    $h(t)=t(\log t-1)$ dengan $h(0)=0$ bersifat konveks ketat pada
    $[0,\infty)$; karena itu, jumlah entropik bersifat konveks ketat dan
    peminimumnya unik.

    Peminimum $P$ tidak dapat mempunyai entri nol. Memang, jika
    $P_{i,j}=0$ untuk sedikitnya satu pasangan $(i,j)$, maka
    $P_t=(1-t)P+tQ$ tetap layak untuk $t\in(0,1)$. Pada setiap entri nol,
    turunan searah satu sisi dari suku entropi memenuhi
    $\frac{\dd}{\dd t}h(tQ_{i,j})=Q_{i,j}\log(tQ_{i,j})
    \rightarrow-\infty$ ketika $t\downarrow0$,
    sedangkan kontribusi biaya linear dan turunan pada entri positif tetap
    hingga. Jadi, untuk $t>0$ yang cukup kecil, nilai objektif pada $P_t$
    lebih kecil daripada pada $P$, suatu kontradiksi. Dengan demikian,
    $P_{i,j}>0$ untuk semua $i,j$, dan kondisi pengali Lagrange untuk kendala
    marginal afin dapat diterapkan.

    Lagrangian masalah ini adalah
    \begin{equation}
        L(P,\lambda,\mu)
        =\inner{C}{P}+\epsilon E(P)
        +\inner{\lambda}{P\1_m-a}
        +\inner{\mu}{P^\top\1_n-b}.
    \end{equation}
    Kondisi optimalitasnya adalah
    \begin{equation}
        \begin{cases}
            0&=C_{i,j}+\epsilon\log(P_{i,j})+\lambda_i+\mu_j,\\
            0&=P\1_m-a,\\
            0&=P^\top\1_n-b.
        \end{cases}
    \end{equation}
    Kondisi pertama menyiratkan
    \begin{equation}
        P_{i,j}
        =\exp\left(-\frac{\lambda_i}{\epsilon}\right)
        \exp\left(-\frac{C_{i,j}}{\epsilon}\right)
        \exp\left(-\frac{\mu_j}{\epsilon}\right).
    \end{equation}
    Dengan $u_i=\exp(-\lambda_i/\epsilon)$ dan
    $v_j=\exp(-\mu_j/\epsilon)$, diperoleh representasi yang dinyatakan.
    Jika pasangan positif $(\tilde u,\tilde v)$ menghasilkan rencana yang sama,
    maka dari $u_iK_{i,j}v_j=\tilde u_iK_{i,j}\tilde v_j$ dan $K_{i,j}>0$
    diperoleh $\tilde u_i/u_i=v_j/\tilde v_j$ untuk setiap $i,j$. Karena ruas
    kiri hanya bergantung pada $i$ dan ruas kanan hanya pada $j$, semua rasio
    itu sama dengan suatu $t>0$; jadi $(\tilde u,\tilde v)=(tu,t^{-1}v)$.
\end{proof}

% H09-S009 | sumber authority/habring/source-v1/optimal_transport.tex baris 257--264
% segment-id: d90.hab.v1.ch09.seg0009
Kendala yang masih harus dipenuhi adalah
$P\1_m=\diag(u)K\diag(v)\1_m=\diag(u)Kv=a$ dan
$P^\top\1_n=\diag(v)K^\top u=b$. Ini memotivasi algoritme
Sinkhorn--Knopp. Pilih $v_0\in\R_{++}^m$, misalnya $v_0=\1_m$, lalu untuk
$k=0,1,2,\dots$ perbarui
\begin{equation}
    \begin{cases}
        u_{k+1}=a\oslash(Kv_k),\\
        v_{k+1}=b\oslash(K^\top u_{k+1}),
    \end{cases}
\end{equation}
dengan $\oslash$ menyatakan pembagian komponen demi komponen. Berdasarkan
positivitas $a$, $b$, dan $K$, semua penyebut dan iterat tetap positif.
Dengan hipotesis pada lema di atas, teorema Sinkhorn--Knopp menyatakan bahwa rencana
$\diag(u_k)K\diag(v_k)$, $k\geq1$, konvergen ke rencana unik pada
masalah~\eqref{ot:eq:disc_entropic}. Vektor penskalaannya sendiri tetap
mempunyai ambiguitas perkalian $(u,v)\mapsto(tu,t^{-1}v)$.


\backmatter
\printbibliography[heading=bibintoc,title={Daftar pustaka}]

\chapter*{Catatan koreksi terhadap sumber}
\markboth{Catatan koreksi terhadap sumber}{Catatan koreksi terhadap sumber}
Terjemahan ini mempertahankan isi sumber kecuali pada koreksi yang dapat ditentukan berikut. Koreksi dilakukan oleh penyusun edisi bahasa Indonesia dan bukan perubahan yang disahkan oleh penulis sumber. Rekaman terstruktur dengan identitas peristiwa O015-HAB-ADV-0084 sampai O015-HAB-ADV-0096 menyertai sumber edisi ini.

\begin{enumerate}
    \item Ruang terukur kini membawa aljabar-$\sigma$ bernama; setiap himpunan yang dipakai pada marginal dan hasil dorong disyaratkan terukur, serta peta transportasi dan fungsi biaya disyaratkan terukur.
    \item Definisi ukuran membedakan ukuran bertanda hingga dari ukuran tak negatif dan melengkapi sifat aljabar-$\sigma$ yang diperlukan untuk kebermaknaan aditivitas terhitung.
    \item Masalah Monge memakai infimum atas peta terukur yang layak, dengan konvensi $\inf\emptyset=+\infty$, karena keberadaan dan ketercapaian peta tidak dijamin.
    \item Contoh ketidaklayakan Monge kini menyatakan $Y\subseteq\R^q$ sebelum menyebut ukuran Lebesgue; catatan kaki menyatakan ketiadaan atom sebagai konsekuensi kontinuitas mutlak, bukan sebagai definisi yang ekuivalen.
    \item Teorema eksistensi kini mensyaratkan ruang Polish Borel. Sketsa bukti memakai keketatan marginal, kekompakan lemah melalui Prokhorov, ketertutupan himpunan kopling, dan semikontinuitas bawah biaya.
    \item Jarak Wasserstein dibangun dari metrik $d$ pada ruang Polish dan dibatasi pada $\Pc_p(X)$, ruang ukuran probabilitas dengan momen ke-$p$ hingga.
    \item Teorema dualitas Kantorovich kini menyatakan ruang Polish Borel, biaya semikontinu bawah tak negatif, serta kelas potensial $C_b(X)\times C_b(Y)$ dan kendala titik demi titik.
    \item Sketsa dualitas membenarkan pemisahan marginal dengan fungsi kontinu terbatas, memperbaiki domain yang salah $\Mc_+(\alpha,\beta)$ menjadi $\Mc_+(X\times Y)$, dan mengidentifikasi secara terbuka pertukaran infimum--supremum sebagai teorema dualitas kuat, bukan akibat manipulasi formal.
    \item Bobot diskret dinyatakan sebagai vektor $a\in\Delta_n$ dan $b\in\Delta_m$; vektor satuan pada kendala marginal diberi dimensi yang benar, transpos dipakai konsisten, dan pemisah kendala yang hilang dipulihkan.
    \item Entropi didefinisikan pada ortan tak negatif dengan $0\log0=0$ dan bernilai $+\infty$ di luarnya; $\epsilon>0$ serta keberhinggaan semua entri matriks biaya dinyatakan sebelum masalah regularisasi.
    \item Lema faktorisasi mensyaratkan marginal positif. Eksistensi dan keunikan dibuktikan dari kekompakan politop dan kekonveksan ketat; perturbasi menuju $ab^\top$ membuktikan positifnya setiap entri sebelum kondisi pengali Lagrange digunakan. Ambiguitas penskalaan juga dinyatakan.
    \item Iterasi Sinkhorn diberi inisialisasi positif dan indeks yang lengkap. Positivitas penyebut serta konvergensi rencana penskalaan ke rencana entropik unik dinyatakan, sambil mempertahankan ambiguitas perkalian pasangan penskalaan.
    \item Entri daftar pustaka Villani memakai pengarang tunggal C\'edric Villani sesuai metadata penerbit; penanda ``and others'' yang keliru pada berkas sumber tidak lagi menghasilkan nama tampilan ``andothers''.
\end{enumerate}

\chapter*{Atribusi, lisensi, dan nondukungan}
Karya sumber: Andreas Habring, \emph{Lecture Notes: Convex Optimization}, arXiv:2607.11664v1, 13 Juli 2026, \url{https://arxiv.org/abs/2607.11664v1}. Karya sumber tersedia berdasarkan Creative Commons Attribution 4.0 International (CC BY 4.0), \url{https://creativecommons.org/licenses/by/4.0/}.

Edisi bahasa Indonesia ini merupakan karya turunan. Hak cipta materi sumber tetap pada pemegang haknya. Anda boleh membagikan dan mengadaptasinya dengan atribusi yang sesuai, tautan lisensi, dan penandaan perubahan sebagaimana disyaratkan CC BY 4.0.

Andreas Habring dan TU Graz tidak menyusun, memeriksa, menyetujui, atau mendukung terjemahan ini. Penyebutan nama mereka semata-mata untuk atribusi karya sumber.

\end{document}
